% Teaching template for 深基课 Chapter 08 (paper writing).
% Example method name: Dual-Path Calibration Network (DPCNet).
% REPLACE the example task/modules/numbers with the student's real project.
% Compile: pdflatex main && bibtex main && pdflatex main && pdflatex main
\documentclass[11pt,a4paper]{article}
\usepackage[margin=1in]{geometry}
\usepackage{times}
\usepackage{amsmath,amssymb,amsfonts}
\usepackage{graphicx}
\usepackage{booktabs}
\usepackage{algorithm}
\usepackage{algpseudocode}
\usepackage{hyperref}
\usepackage{xcolor}

\title{Dual-Path Calibration Network for Robust Visual Recognition\\
\large (Teaching Template --- Replace with Your Title)}
\author{Student Name\quad Mentor Name\\
\small Deep Foundations Course (Shenji) Writing Tutorial}
\date{}

\begin{document}
\maketitle

\begin{abstract}
Visual recognition underpins a wide range of perception systems, yet models still struggle when local evidence and global context are poorly calibrated.
Existing convolutional and transformer pipelines often emphasize either fine-grained textures or long-range dependencies, leaving a persistent gap in jointly stabilizing discriminative cues.
We present the Dual-Path Calibration Network (\textbf{DPCNet}), which couples a local evidence refinement path with a global consistency calibration path under a unified training objective.
The design directly targets two modeling challenges derived from that gap: unreliable local cue selection and weak cross-region consistency.
Experiments on standard classification benchmarks show that DPCNet improves over strong convolutional, transformer, and hybrid baselines, while ablations, sensitivity studies, and visualizations support the contribution of each path.
The template text is instructional; students must replace claims and numbers with results from their own runs.
\end{abstract}

\section{Introduction}
\label{sec:intro}

Deep visual recognition has become a foundational capability for image understanding, retrieval, and downstream perception applications~\cite{deng2009imagenet,krizhevsky2012alexnet,he2016resnet}.
With large-scale supervision and modern architectures, accuracy on canonical benchmarks has improved rapidly, and recognition backbones are routinely transferred to detection and segmentation~\cite{ren2015fasterrcnn,ronneberger2015unet,kirillov2023sam}.
Despite this progress, practical systems still face brittle decisions when textures are ambiguous, objects are partially occluded, or long-range context conflicts with local appearance.
These failure modes matter because recognition errors propagate to any pipeline that treats the classifier as a perception front-end.

Classical convolutional networks provide strong local inductive biases and efficient hierarchical features~\cite{simonyan2015vgg,szegedy2015googlenet,huang2017densenet,xie2017resnext,tan2019efficientnet}, while attention modules further reweight channels or spatial responses~\cite{hu2018senet,woo2018cbam,wang2020eca}.
Vision Transformers and their hierarchical variants model long-range dependencies via self-attention on patches or windows~\cite{dosovitskiy2021vit,liu2021swin,touvron2021deit,wang2021pvt,liu2022swinv2}.
Hybrid and modernized ConvNet designs continue to narrow the gap between convolution and attention~\cite{liu2022convnext,woo2023convnextv2,li2022mvitv2}.
Nevertheless, many pipelines still optimize a single dominant representation route: either local operators dominate and global conflicts remain unresolved, or global attention dominates and locally decisive evidence is diluted.
The core modeling limitation is therefore not the absence of either local or global operators, but the lack of an explicit dual-path calibration mechanism that stabilizes both evidence types under one objective.

This limitation raises two coupled scientific challenges.
First, \emph{local evidence selection} remains unreliable: texture-level responses can be noisy or spuriously correlated, so the model may amplify non-causal cues even when a backbone is strong~\cite{selvaraju2017gradcam}.
Second, \emph{global consistency calibration} remains weak: aggregating long-range context without an explicit consistency constraint can mix incompatible regional semantics, especially under clutter~\cite{vaswani2017attention,dosovitskiy2021vit}.
Existing single-route designs typically mitigate one challenge more than the other, which leaves residual errors that look like either over-localization or over-smoothing.

To address both challenges, we propose the Dual-Path Calibration Network (\textbf{DPCNet}).
DPCNet maintains a local refinement path that reweights discriminative local responses and a global calibration path that enforces cross-region consistency, then fuses the two calibrated streams for classification.
The local path is designed for Challenge~1, while the global path is designed for Challenge~2; their joint objective encourages complementary rather than redundant corrections.
Self-supervised and transfer-oriented representation studies motivate the value of calibrated features~\cite{chen2020simclr,he2020moco,grill2020byol,caron2021dino,he2022mae,radford2021clip}, yet our focus is a supervised dual-path calibration architecture rather than a new pre-training paradigm.

Our contributions are threefold.
(i)~We formulate recognition failures around a single gap---missing dual-path calibration---and derive two explicit challenges on local evidence selection and global consistency.
(ii)~We instantiate DPCNet with paired local and global calibration modules, consistent notation, and a complete training procedure.
(iii)~We provide a full experimental protocol template covering comparisons, ablations, sensitivity, visualization, and error analysis for students to fill with verified runs.

% \begin{figure}[t]
% \centering
% % \includegraphics[width=0.85\linewidth]{fig_motivation_example.png}
% \caption{Illustration of the core modeling limitation of existing methods.}
% \label{fig:motivation_example}
% \end{figure}

\section{Related Work}
\label{sec:rw}

\subsection{Convolutional Recognition and Local Attention}
Krizhevsky et al.~\cite{krizhevsky2012alexnet} established deep CNNs as a practical recognition stack, and subsequent residual, dense, and aggregated designs improved optimization and feature reuse~\cite{he2016resnet,huang2017densenet,xie2017resnext}.
Tan and Le~\cite{tan2019efficientnet} further studied compound scaling for accuracy--efficiency trade-offs, while mobile architectures targeted resource-constrained deployment~\cite{howard2017mobilenets,sandler2018mobilenetv2}.
Channel and spatial attention blocks such as SE, CBAM, and ECA reweight local responses with limited overhead~\cite{hu2018senet,woo2018cbam,wang2020eca}.
These methods are strong at encoding neighborhood structure, yet their receptive-field growth is primarily hierarchical; when distant regions disagree, local attention alone may not impose an explicit consistency correction.
Liu et al.~\cite{liu2022convnext} and Woo et al.~\cite{woo2023convnextv2} modernize ConvNets with transformer-inspired training recipes, improving competitiveness, but still center a convolution-dominated route rather than dual-path calibration between local evidence and global consistency.

\subsection{Transformer-Based and Hybrid Global Modeling}
Dosovitskiy et al.~\cite{dosovitskiy2021vit} showed that patch-based Transformers can match or surpass CNNs when pre-trained at scale, and DeiT improved data efficiency via distillation~\cite{touvron2021deit}.
Hierarchical Transformers such as Swin and PVT introduce multi-scale tokens and windowed attention for denser prediction friendliness~\cite{liu2021swin,wang2021pvt,liu2022swinv2}.
Alternative token mixing designs, including MLP-Mixer and Tokens-to-Token ViT, explore different global aggregation operators~\cite{tolstikhin2021mlpmixer,yuan2021t2t,chu2021twins}.
Self-supervised objectives further strengthen transferable features~\cite{chen2020simclr,he2020moco,chen2021mocov3,caron2021dino,zhou2022ibot,he2022mae,bao2022beit,oquab2024dinov2}.
These lines excel at long-range interaction, yet global mixing can over-smooth decisive local cues if no dedicated local refinement path is co-trained for calibration.
In contrast to replacing one backbone family with another, DPCNet keeps an explicit local--global dual path and focuses on their joint calibration for recognition.

\section{Method}
\label{sec:method}

\subsection{Problem Formulation}
\label{sec:formulation}
Let an input image be $\mathbf{X}\in\mathbb{R}^{H\times W\times C}$ with label $y\in\{1,\ldots,K\}$.
A backbone encoder $f_{\theta}$ maps $\mathbf{X}$ to an intermediate feature map $\mathbf{F}\in\mathbb{R}^{h\times w\times d}$:
\begin{equation}
\mathbf{F}=f_{\theta}(\mathbf{X}).
\label{eq:backbone}
\end{equation}
Our goal is to learn a calibrated representation $\mathbf{z}\in\mathbb{R}^{d}$ and a classifier $g_{\phi}$ that minimizes the expected risk
\begin{equation}
\min_{\theta,\phi,\psi}\;\mathbb{E}_{(\mathbf{X},y)}\big[\mathcal{L}_{\mathrm{cls}}(g_{\phi}(\mathbf{z}),y)+\lambda\mathcal{L}_{\mathrm{cal}}(\mathbf{F};\psi)\big],
\label{eq:obj}
\end{equation}
where $\psi$ collects calibration-module parameters, $\mathcal{L}_{\mathrm{cls}}$ is cross-entropy, and $\mathcal{L}_{\mathrm{cal}}$ regularizes dual-path consistency.
All vectors are column vectors unless noted; $\|\cdot\|_2$ denotes the Euclidean norm.

\subsection{Model Overview}
\label{sec:overview}
DPCNet processes $\mathbf{F}$ through a Local Evidence Refinement (LER) module and a Global Consistency Calibration (GCC) module, then fuses their outputs.
LER addresses unreliable local cue selection; GCC addresses weak cross-region consistency.
The overall data flow is: encode $\rightarrow$ refine locally $\rightarrow$ calibrate globally $\rightarrow$ fuse $\rightarrow$ classify.
% Figure~\ref{fig:model_arch} illustrates the architecture.

\subsection{Local Evidence Refinement}
\label{sec:ler}
Given $\mathbf{F}$, we compute a local importance map $\mathbf{A}_{\ell}\in\mathbb{R}^{h\times w}$ and a refined field $\mathbf{F}_{\ell}$:
\begin{equation}
\mathbf{A}_{\ell}=\sigma\big(\mathrm{Conv}_{1\times1}(\mathbf{F})\big),
\label{eq:local_attn}
\end{equation}
\begin{equation}
\mathbf{F}_{\ell}=\mathbf{F}\odot\mathrm{Expand}(\mathbf{A}_{\ell}),
\label{eq:local_refine}
\end{equation}
where $\sigma$ is the sigmoid, $\odot$ is element-wise multiplication, and $\mathrm{Expand}$ broadcasts channel-wise.
A pooled local descriptor is
\begin{equation}
\mathbf{z}_{\ell}=\mathrm{GAP}(\mathbf{F}_{\ell})\in\mathbb{R}^{d},
\label{eq:local_pool}
\end{equation}
with $\mathrm{GAP}$ denoting global average pooling.
Intuitively, Eq.~\eqref{eq:local_attn}--\eqref{eq:local_pool} suppress non-informative locations before aggregation.

\subsection{Global Consistency Calibration}
\label{sec:gcc}
We flatten $\mathbf{F}$ into tokens $\mathbf{T}\in\mathbb{R}^{n\times d}$ with $n=hw$ and apply multi-head self-attention~\cite{vaswani2017attention}:
\begin{equation}
\mathbf{T}'=\mathrm{MHSA}(\mathbf{T})+\mathbf{T},
\label{eq:mhsa}
\end{equation}
\begin{equation}
\mathbf{T}_{g}=\mathrm{FFN}(\mathbf{T}')+\mathbf{T}'.
\label{eq:ffn}
\end{equation}
A consistency regularizer encourages agreement among high-confidence tokens:
\begin{equation}
\mathcal{L}_{\mathrm{cons}}=\frac{1}{|\mathcal{P}|}\sum_{(i,j)\in\mathcal{P}}\big\|\hat{\mathbf{t}}_{i}-\hat{\mathbf{t}}_{j}\big\|_2^{2},
\label{eq:cons}
\end{equation}
where $\hat{\mathbf{t}}_{i}$ is an $\ell_2$-normalized token and $\mathcal{P}$ indexes pairs above an attention threshold.
The global descriptor is
\begin{equation}
\mathbf{z}_{g}=\mathrm{GAP}(\mathrm{Reshape}(\mathbf{T}_{g})).
\label{eq:global_pool}
\end{equation}

\subsection{Fusion and Prediction}
\label{sec:fuse}
We fuse the two descriptors with a gated sum:
\begin{equation}
\boldsymbol{\alpha}=\mathrm{softmax}(\mathbf{W}_{g}[\mathbf{z}_{\ell};\mathbf{z}_{g}]),
\label{eq:gate}
\end{equation}
\begin{equation}
\mathbf{z}=\alpha_{1}\mathbf{z}_{\ell}+\alpha_{2}\mathbf{z}_{g},
\label{eq:fuse}
\end{equation}
\begin{equation}
\hat{\mathbf{y}}=\mathrm{softmax}(\mathbf{W}_{c}\mathbf{z}+\mathbf{b}_{c}).
\label{eq:pred}
\end{equation}
The classification loss is
\begin{equation}
\mathcal{L}_{\mathrm{cls}}=-\sum_{k=1}^{K}y_{k}\log\hat{y}_{k},
\label{eq:ce}
\end{equation}
and the full objective instantiates Eq.~\eqref{eq:obj} as
\begin{equation}
\mathcal{L}=\mathcal{L}_{\mathrm{cls}}+\lambda\mathcal{L}_{\mathrm{cons}}.
\label{eq:total}
\end{equation}

\subsection{Training Procedure}
\label{sec:train}
\begin{algorithm}[t]
\caption{Training procedure of DPCNet}
\label{alg:train}
\begin{algorithmic}[1]
\Require Dataset $\mathcal{D}$, backbone $f_{\theta}$, modules $\psi$, classifier $g_{\phi}$, learning rate $\eta$, weight $\lambda$
\Ensure Optimized parameters $\{\theta,\psi,\phi\}$
\For{epoch $=1$ to $E$}
  \For{each minibatch $(\mathbf{X},y)\sim\mathcal{D}$}
    \State $\mathbf{F}\leftarrow f_{\theta}(\mathbf{X})$
    \State Compute $\mathbf{z}_{\ell},\mathbf{z}_{g}$ via LER and GCC
    \State Fuse $\mathbf{z}$ and predict $\hat{\mathbf{y}}$
    \State $\mathcal{L}\leftarrow\mathcal{L}_{\mathrm{cls}}+\lambda\mathcal{L}_{\mathrm{cons}}$
    \State Update $\{\theta,\psi,\phi\}$ with AdamW~\cite{loshchilov2019adamw}
  \EndFor
\EndFor
\end{algorithmic}
\end{algorithm}

\subsection{Complexity Analysis}
\label{sec:complexity}
Let $n=hw$ be the token count and $d$ the channel width.
LER is dominated by $1\times1$ convolution and pooling, i.e., $\mathcal{O}(nd)$ time per image.
GCC attention is $\mathcal{O}(n^{2}d)$ in the dense case (or $\mathcal{O}(nwd)$ with window size $w$ if windowed attention is used~\cite{liu2021swin}).
Training stores activations for backpropagation, yielding $\mathcal{O}(nd+Ed)$ peak feature memory aside from optimizer states.
Relative to a pure CNN baseline, the extra cost is mainly the global path; relative to a full Transformer baseline, LER adds only linear overhead while providing an explicit local calibration branch.

\section{Experiments}
\label{sec:exp}

\subsection{Experimental Setup}
\label{sec:setup}
\textbf{Datasets.}
We evaluate on ImageNet-style classification settings and mid-scale benchmarks commonly used for backbone comparison~\cite{deng2009imagenet}.
Students should state split sizes, preprocessing, and licenses for every dataset they actually run.

\textbf{Baselines.}
We compare three families: (i)~CNN baselines (ResNet, DenseNet, EfficientNet, ConvNeXt)~\cite{he2016resnet,huang2017densenet,tan2019efficientnet,liu2022convnext};
(ii)~Transformer baselines (ViT, DeiT, Swin)~\cite{dosovitskiy2021vit,touvron2021deit,liu2021swin};
(iii)~attention-augmented CNNs (SE, CBAM)~\cite{hu2018senet,woo2018cbam}.
All methods must be reproduced under the same split and training recipe for fair comparison.

\textbf{Metrics.}
We report Top-1 / Top-5 accuracy (or task-specific F1/mAP).
Define each metric equation in the student paper when non-standard.

\textbf{Implementation.}
Unless noted, we use AdamW~\cite{loshchilov2019adamw}, cosine decay, and standard augmentations such as Mixup/CutMix/RandAugment when applicable~\cite{zhang2018mixup,yun2019cutmix,cubuk2020randaugment}.
List GPU type, epoch count, batch size, and seeds in the final manuscript.

\subsection{Performance Comparison}
\label{sec:comparison}
% TEMPLATE PARAGRAPHS --- replace bracketed text with verified numbers from your runs.
DPCNet attains the best overall accuracy among reproduced baselines under the shared protocol.
On the primary benchmark, it reaches [TOP1]\% Top-1, exceeding the strongest CNN baseline by [D1] points and the strongest Transformer baseline by [D2] points, which indicates that joint local--global calibration is more effective than strengthening either route alone.

Strong CNN baselines outperform lightweight mobile CNNs because residual multi-scale features preserve richer local evidence for LER to reweight, whereas aggressively compressed mobile towers discard cues before calibration.
Quantitatively, ResNet/ConvNeXt-style models sit [G1] points above MobileNet-style counterparts on Top-1 in our reproduction table, consistent with capacity and feature fidelity differences reported in prior scaling studies~\cite{tan2019efficientnet,howard2017mobilenets}.

Transformer baselines outperform classical attention-CNN hybrids on long-range dominated cases, yet can lag when local textures decide the label.
In our table, Swin/DeiT exceed SE/CBAM by [G2] points on average Top-1, while DPCNet further improves by coupling GCC with an explicit LER path rather than relying on global mixing alone~\cite{liu2021swin,touvron2021deit,woo2018cbam}.

Overall, the ranking supports a mechanistic reading: methods weak on local fidelity underperform when fine textures matter; methods weak on global consistency underperform under clutter; DPCNet reduces both error modes in the same model.
% Insert Table~\ref{tab:comparison} with >=8 reproduced methods.

\subsection{Ablation Study}
\label{sec:ablation}
To validate each component, we evaluate: full DPCNet; w/o LER; w/o GCC; w/o consistency loss $\mathcal{L}_{\mathrm{cons}}$; and a fusion-replacement variant that concatenates $\mathbf{z}_{\ell},\mathbf{z}_{g}$ without gating.
Removing LER primarily hurts fine-grained categories, confirming Challenge~1.
Removing GCC increases confusion among context-heavy classes, confirming Challenge~2.
Removing $\mathcal{L}_{\mathrm{cons}}$ preserves architecture but weakens consistency, yielding a smaller yet consistent drop.
[Replace with table numbers and paragraph-length analysis tied to mechanisms.]

\subsection{Parameter Analysis}
\label{sec:param}
We study the consistency weight $\lambda$ and the LER bottleneck width.
Performance peaks at moderate $\lambda$: too small under-regularizes global tokens, too large over-smooths class boundaries.
Width sweeps show diminishing returns beyond a compact setting, suggesting calibration quality matters more than merely widening the local gate.
[Plot curves from real sweeps; explain with the same dual-challenge language.]

\subsection{Visualization Analysis}
\label{sec:vis}
To examine whether DPCNet improves separability, we compare baseline versus DPCNet features with t-SNE/UMAP~\cite{vandermaaten2008tsne,mcinnes2018umap} on a fixed probe set and seed.
We further visualize LER maps and Grad-CAM responses~\cite{selvaraju2017gradcam} to check whether local refinement focuses on object support rather than background.
Cleaner clusters and more object-aligned highlights provide qualitative support for the two calibration paths.
[Describe actual figures after generation.]

\subsection{Case and Error Analysis}
\label{sec:error}
We present cases where the baseline fails but DPCNet succeeds, typically when local texture is weak yet global structure is consistent, or the reverse.
Residual errors concentrate on fine-grained pairs with nearly identical local patterns, indicating that dual-path calibration mitigates but does not eliminate intrinsic ambiguity.
These observations bound the method's applicability and motivate future disambiguation modules.

\section{Conclusion}
\label{sec:concl}

\subsection{Summary}
This work revisited visual recognition through the lens of missing dual-path calibration between local evidence and global consistency.
We derived two challenges---unreliable local cue selection and weak cross-region consistency---and introduced DPCNet to address them with LER and GCC under a unified objective.
Template experiments outline how comparisons, ablations, and analyses should substantiate that the full model surpasses strong single-route baselines and its own degraded variants.

\subsection{Limitations}
DPCNet assumes that useful local and global cues coexist in the same feature map and that pairwise token consistency is a meaningful surrogate for semantic agreement.
When categories are defined by compositional rules beyond appearance similarity, consistency regularization may align the wrong tokens.
The gated fusion also assumes that a static mixture of two descriptors is sufficient; tasks requiring conditional routing across many experts may expose this restriction.

\subsection{Future Work}
Corresponding to the first limitation, future work may replace appearance-based consistency with structure-aware constraints that incorporate part--whole relations.
Corresponding to the second limitation, extending gated fusion toward conditional computation or mixture-of-experts routing is a concrete next step.
A further direction is to test whether the same dual-path calibration transfers to dense prediction heads without changing the challenge formulation.

\bibliographystyle{plain}
\bibliography{refs}

\end{document}
